\section{Purpose, normalization, and evidence}
Put $\theta=\beta(4)=\sum_{k\geq0}(-1)^k/(2k+1)^4$.
The companion paper~\cite{Main} proves two primitive-value obstructions
for positive determinant families. The present supplement supplies the
broader mathematical record behind that selection. It organizes the
investigation by construction, rather than by the chronology of eighteen
research rounds. Stable record identifiers, such as S03 or P07, refer to
the source register in Appendix~\ref{app:register} and the repository's
57-entry inventory. They are navigation labels, not theorem numbers.

The motivating determinant work of OpenAI~\cite[Section~1.2]{OpenAI}
separates nonvanishing, real size and arithmetic normalization.
Zudilin~\cite{Zudilin} gives the broader positive-moment determinant
framework. The distinction between real decay and primitive arithmetic
remained decisive in this investigation. Carreira's work~\cite{Carreira}
provides directly related $\beta(4)$ constructions and tests; the general
positive-weight and rationally modified moment frameworks also appear in
Pr\'evost~\cite{Prevost} and Krattenthaler~\cite{Krattenthaler}.
These are sources of context and specific attributed constructions,
not inputs asserting a global irrationality theorem.

For a nonzero $Q\in\Q[X]$, write $\Prim(Q)=\kappa Q$, where
$\kappa>0$ and the resulting integer coefficients have gcd one. Then
\begin{equation}\label{eq:primitive}
 v_p(\kappa)=-\min_j v_p([X^j]Q).
\end{equation}
This identity follows because the minimum valuation of a primitive
integer coefficient list is zero at every prime. Multiplying $Q$ by a
nonzero rational scalar preserves its primitive representative up to
sign. A lower bound on a convenient clearing integer is therefore not
automatically a lower bound on $\kappa$. For example, $5$ clears
$(6X+12)/5$, whose primitive multiplier is $5/6$.

If $\theta=a/b$ and $Z\in\Z[X]$ has degree at most $d$, then
$b^d Z(\theta)\in\Z$. Thus nonzero values satisfying
$|Z_n(\theta)|^{1/d_n}\to0$ would prove irrationality. Merely having
$Z_n(\theta)\to0$ with unbounded degrees is insufficient: $X^n$ at
$1/2$ is a counterexample. For linear forms this reduces to the usual
nonzero integer-form criterion.

Four conclusions must be kept separate. A \emph{primitive obstruction}
bounds actual primitive values. An \emph{estimate limitation} says only
that a specified bound or normalization is inadequate. A \emph{structural
obstruction} disproves a particular proposed implication or transformation.
\emph{Finite evidence} concerns only the enumerated instances. The method
record preserves these distinctions, including where a later correction
weakened an earlier interpretation. No statement about all possible
irrationality arguments follows from these results.

Standard background includes the prime number theorem
$\vartheta(x)=\sum_{p\leq x}\log p=x+o(x)$, elementary complex integration,
Fourier convergence and the stated elementary polynomial identities.
The completed Mellin application has an additional, explicitly identified
polylogarithmic input in Section~\ref{sec:mellin}. The source register
distinguishes proofs rewritten here from results retained as archived
claims. No new independent audit of this assembled supplement is asserted.

\section{Positive polynomial sources: a scope boundary}\label{sec:cone}
Define
\begin{equation}\label{eq:scalar}
 \dd\nu(t)=\frac{(-\log t)^3}{96\sqrt t}\dd t,
 \qquad I(P)=\int_0^1\frac{P(t)}{1+t}\dd\nu(t).
\end{equation}
Here $\nu_j=\int t^j\dd\nu=(2j+1)^{-4}$ and $I(1)=\theta$.
For a rational polynomial define its \emph{formal pair} by
\begin{equation}\label{eq:formal-pair}
 A(P)=P(-1),\qquad
 B(P)=\int_0^1\frac{P(t)-P(-1)}{1+t}\dd\nu(t).
\end{equation}
The quotient is a polynomial, so $B(P)\in\Q$ and
$I(P)=A(P)\theta+B(P)$. This defines the coefficients without assuming
that $1$ and $\theta$ are linearly independent.

\begin{theorem}[The unrestricted positive cone; S01]\label{thm:cone}
The formal-pair image of $P\in\Q[t]$ nonnegative on $[0,1]$ is
\[
 \{(0,0)\}\ \cup\ \{(A,B)\in\Q^2:A\theta+B>0\}.
\]
The same positive primitive integer pairs are obtained if the source
polynomials are required to have integer coefficients.
\end{theorem}
\begin{proof}
Nonzero nonnegative polynomials have strictly positive integrals. Conversely
let $L=A\theta+B>0$. If $A=0$, use $P=B(1+t)$, since $\nu_0=1$.
For $A>0$ use $P_n=A((1-t)/2)^{2n}$; for $A<0$ use
$P_n=(-A)t((1-t)/2)^{2n}$. Both have $A(P_n)=A$, are nonnegative,
and satisfy $I(P_n)\leq |A|\theta4^{-n}$. Choose $n$ making this less
than $L$. Then
\[
 P=P_n+(B-B(P_n))(1+t)
\]
is nonnegative and has exactly the desired pair; the added coefficient
is rational and equals $L-I(P_n)>0$. Clearing the rational coefficients
of $P$ multiplies its pair by a positive integer. Primitive normalization
returns the original primitive pair.
\end{proof}

This construction encodes an already specified approximation; it does
not produce a new one. In particular, shrinking positive primitive
integer forms in this unrestricted cone exist if and only if $\theta$
is irrational. Necessity follows from the rational lower bound $1/b$;
sufficiency follows from rational approximation to an irrational real
number and Theorem~\ref{thm:cone}. A negative theorem for all such sources
would therefore be a rationality theorem, not a consequence of failures
of a few prescribed weights.

\begin{proposition}[Cost in degree; S01]\label{prop:degree}
Let $\rho=3+2\sqrt2$. For $A\ne0$ and $L=A\theta+B>0$, the least
degree $d(A,B)$ of a nonnegative rational witness satisfies
\[
 d(A,B)=\frac{\log(|A|/L)}{\log\rho}
             +O\bigl(\log\log(|A|/L)\bigr)
 \qquad (|A|/L\longrightarrow\infty).
\]
\end{proposition}
\begin{proof}
Let $q_n(t)=T_n(2t-1)/T_n(-3)$, where $T_n$ is the Chebyshev
polynomial. Then $q_n(-1)=1$ and
$\|q_n\|_{[0,1]}\leq 2\rho^{-n}$. Replace $P_n$ in the preceding
proof by $Aq_n^2$ or $(-A)tq_n^2$. Its integral is at most
$4|A|\rho^{-2n}$ and its degree at most $2n+1$. This gives the upper
bound, with an absolute additive constant.

For a degree-$D$ nonnegative $P$ set $M=\max_{[0,1]}P$.
Its Chebyshev expansion has constant coefficient at most $M$ in absolute
value and other coefficients at most $2M$, by the cosine coefficient
integrals. The identity $T_k(\cos u)=\cos(ku)$ gives
$\|T_k'\|_{[-1,1]}\leq k^2$, hence
$\|P'\|_{[0,1]}\leq4MD^3$. On an interval of length
$\ell=1/(8D^3)$ adjoining a maximizer, $P\geq M/2$.
On its central half, $1-t\geq\ell/4$, and
$\dd\nu/(1+t)\geq(1-t)^3\dd t/192$. Thus
\[
 I(P)\geq\frac{M}{3\cdot2^{26}D^{12}}.
\]
Interpolation at the $D+1$ Chebyshev extrema gives
$|P(-1)|\leq M T_D(3)\leq M\rho^D$: the exterior Lagrange
coefficients alternate in sign, and their absolute sum is $T_D(3)$.
Therefore
\begin{equation}\label{eq:cone-lower}
 I(P)\geq\frac{|P(-1)|}{3\cdot2^{26}D^{12}\rho^D}.
\end{equation}
Combining this with the upper construction proves the assertion.
\end{proof}

The archived bounded-degree extremizers and a degree-16 witness for
$5719097-5783031\theta>0$ are finite evidence
\record{S02}. Their role is to check scope: unfavorable boundary
examples do not exclude favorable interior forms. The complete
bounded-degree boundary formulas are retained in the source record;
they are not needed for either proof above.
